
\documentclass{article}
\usepackage{colortbl}
\usepackage{makecell}
\usepackage{multirow}
\usepackage{supertabular}

\begin{document}

\newcounter{utterance}

\centering \large Interaction Transcript for game `codenames', experiment `association\_difficult', episode 3 with gemma{-}4{-}26b{-}a4b{-}it.
\vspace{24pt}

{ \footnotesize  \setcounter{utterance}{1}
\setlength{\tabcolsep}{0pt}
\begin{supertabular}{c@{$\;$}|p{.15\linewidth}@{}p{.15\linewidth}p{.15\linewidth}p{.15\linewidth}p{.15\linewidth}p{.15\linewidth}}
    \# & $\;$A & \multicolumn{4}{c}{Game Master} & $\;\:$B\\
    \hline

    \theutterance \stepcounter{utterance}  
    & & & \multicolumn{2}{p{0.3\linewidth}}{
        \cellcolor[rgb]{0.9,0.9,0.9}{
            \makecell[{{p{\linewidth}}}]{
                \texttt{\tiny{[GM$|$GM]}}
                {'hidden': {'team': ['robin', 'palec', 'stonoga', 'malarija', 'zajemalka', 'cimbala', 'pivo', 'bež', 'poljska'], 'innocent': ['molj', 'norveška', 'centimeter', 'metlica', 'hrošč', 'oriole', 'viski'], 'opponent': ['rak', 'milja', 'rog', 'mehika', 'kardinal', 'črv', 'sladkorna bolezen', 'teal'], 'assassin': ['rum']}, 'revealed': {'team': {'team': [], 'innocent': [], 'opponent': [], 'assassin': []}, 'opponent': {'team': [], 'innocent': [], 'opponent': [], 'assassin': []}}}
            }
        }
    }
    & & \\ \\

    \theutterance \stepcounter{utterance}  
    & & \multicolumn{4}{p{0.6\linewidth}}{
        \cellcolor[rgb]{0.9,0.9,0.9}{
            \makecell[{{p{\linewidth}}}]{
                \texttt{\tiny{[P1$\langle$GM]}}
                \texttt{Navedite eno enobesedno iztočnico, ki se nanaša na nekatere besede vaše ekipe s spodnjega seznama. Izberete lahko, na koliko in na katere besede se nanaša vaša iztočnica.} \\
\texttt{Cilj je, da vaš partner na podlagi vaših iztočnic ugane vse besede vaše ekipe, preden druga ekipa odkrije vse svoje besede. Če vaš partner ugane besedo "nasprotnika" ali "motečo" besedo, se njegovo ugibanje konča in na vrsto pride nasprotna ekipa. Če vaš partner ugane besedo "ubij", se igra konča in vaša ekipa takoj izgubi.} \\
\texttt{Več besed, na katere se nanaša vaša iztočnica, pomeni hitrejšo zmago v igri, vendar bo vaš partner morda težje uganil asociacije, zato pametno izberite iztočnico. Beseda iztočnica mora biti pomensko povezana s ciljnimi besedami, ne sme biti ena od besed na seznamih ali vsebovati njihovih delov.} \\
\\ 
\texttt{Vedno navedite enobesedno iztočnico in z vejico ločen seznam sorodnih ciljnih besed v naslednji obliki, odgovori naj bodo čim krajši, nikoli ne vključite drugega besedila, kot je zahtevano v tej obliki:} \\
\\ 
\texttt{NAMIG: <BESEDA>} \\
\texttt{CILJI: <CILJI>} \\
\\ 
\texttt{Vaše ekipne besede so: robin, palec, stonoga, malarija, zajemalka, cimbala, pivo, bež, poljska.} \\
\texttt{Besede nasprotnikove ekipe so: rak, milja, rog, mehika, kardinal, črv, sladkorna bolezen, teal.} \\
\texttt{Besede, ki odvračajo pozornost, so: molj, norveška, centimeter, metlica, hrošč, oriole, viski.} \\
\texttt{Besede za ubijanje so: rum.} \\
            }
        }
    }
    & \\ \\

    \theutterance \stepcounter{utterance}  
    & \multicolumn{4}{p{0.6\linewidth}}{
        \cellcolor[rgb]{0.9,0.9,0.9}{
            \makecell[{{p{\linewidth}}}]{
                \texttt{\tiny{[P1$\rangle$GM]}}
                \texttt{NAMIG: narava} \\
\texttt{CILJI: robin, stonoga, malarija, poljska} \\
            }
        }
    }
    & & \\ \\

    \theutterance \stepcounter{utterance}  
    & & & \multicolumn{2}{p{0.3\linewidth}}{
        \cellcolor[rgb]{0.9,0.9,0.9}{
            \makecell[{{p{\linewidth}}}]{
                \texttt{\tiny{[GM$|$GM]}}
                \texttt{narava} \\
            }
        }
    }
    & & \\ \\

    \theutterance \stepcounter{utterance}  
    & & & \multicolumn{2}{p{0.3\linewidth}}{
        \cellcolor[rgb]{0.9,0.9,0.9}{
            \makecell[{{p{\linewidth}}}]{
                \texttt{\tiny{[GM$|$GM]}}
                ['robin', 'stonoga', 'malarija', 'poljska']
            }
        }
    }
    & & \\ \\

    \theutterance \stepcounter{utterance}  
    & & & \multicolumn{2}{p{0.3\linewidth}}{
        \cellcolor[rgb]{0.9,0.9,0.9}{
            \makecell[{{p{\linewidth}}}]{
                \texttt{\tiny{[GM$|$GM]}}
                {'word': 'robin', 'assignment': 'team'}
            }
        }
    }
    & & \\ \\

    \theutterance \stepcounter{utterance}  
    & & & \multicolumn{2}{p{0.3\linewidth}}{
        \cellcolor[rgb]{0.9,0.9,0.9}{
            \makecell[{{p{\linewidth}}}]{
                \texttt{\tiny{[GM$|$GM]}}
                {'word': 'stonoga', 'assignment': 'team'}
            }
        }
    }
    & & \\ \\

    \theutterance \stepcounter{utterance}  
    & & & \multicolumn{2}{p{0.3\linewidth}}{
        \cellcolor[rgb]{0.9,0.9,0.9}{
            \makecell[{{p{\linewidth}}}]{
                \texttt{\tiny{[GM$|$GM]}}
                {'word': 'malarija', 'assignment': 'team'}
            }
        }
    }
    & & \\ \\

    \theutterance \stepcounter{utterance}  
    & & & \multicolumn{2}{p{0.3\linewidth}}{
        \cellcolor[rgb]{0.9,0.9,0.9}{
            \makecell[{{p{\linewidth}}}]{
                \texttt{\tiny{[GM$|$GM]}}
                {'word': 'poljska', 'assignment': 'team'}
            }
        }
    }
    & & \\ \\

    \theutterance \stepcounter{utterance}  
    & & \multicolumn{4}{p{0.6\linewidth}}{
        \cellcolor[rgb]{0.9,0.9,0.9}{
            \makecell[{{p{\linewidth}}}]{
                \texttt{\tiny{[GM$\rangle$P2]}}
                \texttt{Navedite z vejico ločen seznam do 4 besed s spodnjega seznama, ki so najbolj povezane ali najbolj tesno povezane z besedo 'narava'. Seznam ugibanj vedno začnite z UGIBANJE: in v odgovor ne vključite nobenega drugega besedila.} \\
\\ 
\texttt{malarija, norveška, robin, kardinal, viski, sladkorna bolezen, metlica, bež, poljska, rak, oriole, črv, mehika, stonoga, zajemalka, milja, palec, rum, pivo, centimeter, teal, molj, hrošč, cimbala, rog} \\
            }
        }
    }
    & \\ \\

    \theutterance \stepcounter{utterance}  
    & & & \multicolumn{4}{p{0.6\linewidth}}{
        \cellcolor[rgb]{0.9,0.9,0.9}{
            \makecell[{{p{\linewidth}}}]{
                \texttt{\tiny{[GM$\langle$P2]}}
                \texttt{UGIBANJE: črv, stonoga, hrošč, molj} \\
            }
        }
    }
     \\ \\

    \theutterance \stepcounter{utterance}  
    & & & \multicolumn{2}{p{0.3\linewidth}}{
        \cellcolor[rgb]{0.9,0.9,0.9}{
            \makecell[{{p{\linewidth}}}]{
                \texttt{\tiny{[GM$|$GM]}}
                ['črv', 'stonoga', 'hrošč', 'molj']
            }
        }
    }
    & & \\ \\

    \theutterance \stepcounter{utterance}  
    & & & \multicolumn{2}{p{0.3\linewidth}}{
        \cellcolor[rgb]{0.9,0.9,0.9}{
            \makecell[{{p{\linewidth}}}]{
                \texttt{\tiny{[GM$|$GM]}}
                {'word': 'črv', 'assignment': 'opponent'}
            }
        }
    }
    & & \\ \\

    \theutterance \stepcounter{utterance}  
    & & & \multicolumn{2}{p{0.3\linewidth}}{
        \cellcolor[rgb]{0.9,0.9,0.9}{
            \makecell[{{p{\linewidth}}}]{
                \texttt{\tiny{[GM$|$GM]}}
                \texttt{črv} \\
            }
        }
    }
    & & \\ \\

    \theutterance \stepcounter{utterance}  
    & & & \multicolumn{2}{p{0.3\linewidth}}{
        \cellcolor[rgb]{0.9,0.9,0.9}{
            \makecell[{{p{\linewidth}}}]{
                \texttt{\tiny{[GM$|$GM]}}
                {'word': 'milja', 'assignment': 'opponent'}
            }
        }
    }
    & & \\ \\

    \theutterance \stepcounter{utterance}  
    & & & \multicolumn{2}{p{0.3\linewidth}}{
        \cellcolor[rgb]{0.9,0.9,0.9}{
            \makecell[{{p{\linewidth}}}]{
                \texttt{\tiny{[GM$|$GM]}}
                {'hidden': {'team': ['robin', 'palec', 'stonoga', 'malarija', 'zajemalka', 'cimbala', 'pivo', 'bež', 'poljska'], 'innocent': ['molj', 'norveška', 'centimeter', 'metlica', 'hrošč', 'oriole', 'viski'], 'opponent': ['rak', 'rog', 'mehika', 'kardinal', 'sladkorna bolezen', 'teal'], 'assassin': ['rum']}, 'revealed': {'team': {'team': [], 'innocent': [], 'opponent': ['črv'], 'assassin': []}, 'opponent': {'team': [], 'innocent': [], 'opponent': ['milja'], 'assassin': []}}}
            }
        }
    }
    & & \\ \\

    \theutterance \stepcounter{utterance}  
    & & \multicolumn{4}{p{0.6\linewidth}}{
        \cellcolor[rgb]{0.9,0.9,0.9}{
            \makecell[{{p{\linewidth}}}]{
                \texttt{\tiny{[P1$\langle$GM]}}
                \texttt{Beseda črv je bila uganjena, vendar je nasprotnik beseda. Potezo vašega soigralca se je tam končala.} \\
\texttt{Sedaj podajte še en namig, povezan z nekaterimi vašimi preostalimi besedami ekipe, in seznam povezanih ciljnih besed. Ne pozabite začeti namiga z 'NAMIG: ', nato dodati novo vrstico in začeti seznam ciljnih besed, ločenih z vejico, z 'CILJI: '. Opomba: nekatere besede so bile s seznamov odstranjene v primerjavi s prejšnjimi zahtevami.} \\
\\ 
\texttt{Vaše preostale besede ekipe so: robin, palec, stonoga, malarija, zajemalka, cimbala, pivo, bež, poljska.} \\
\texttt{Preostale besede nasprotnika so: rak, milja, rog, mehika, kardinal, sladkorna bolezen, teal.} \\
\texttt{Preostale besede, ki motijo, so: molj, norveška, centimeter, metlica, hrošč, oriole, viski.} \\
\texttt{Preostale smrtonosne besede so: rum.} \\
            }
        }
    }
    & \\ \\

    \theutterance \stepcounter{utterance}  
    & \multicolumn{4}{p{0.6\linewidth}}{
        \cellcolor[rgb]{0.9,0.9,0.9}{
            \makecell[{{p{\linewidth}}}]{
                \texttt{\tiny{[P1$\rangle$GM]}}
                \texttt{NAMIG: glasba} \\
\texttt{CILJI: cimbala, pivo} \\
            }
        }
    }
    & & \\ \\

    \theutterance \stepcounter{utterance}  
    & & & \multicolumn{2}{p{0.3\linewidth}}{
        \cellcolor[rgb]{0.9,0.9,0.9}{
            \makecell[{{p{\linewidth}}}]{
                \texttt{\tiny{[GM$|$GM]}}
                \texttt{glasba} \\
            }
        }
    }
    & & \\ \\

    \theutterance \stepcounter{utterance}  
    & & & \multicolumn{2}{p{0.3\linewidth}}{
        \cellcolor[rgb]{0.9,0.9,0.9}{
            \makecell[{{p{\linewidth}}}]{
                \texttt{\tiny{[GM$|$GM]}}
                ['cimbala', 'pivo']
            }
        }
    }
    & & \\ \\

    \theutterance \stepcounter{utterance}  
    & & & \multicolumn{2}{p{0.3\linewidth}}{
        \cellcolor[rgb]{0.9,0.9,0.9}{
            \makecell[{{p{\linewidth}}}]{
                \texttt{\tiny{[GM$|$GM]}}
                {'word': 'cimbala', 'assignment': 'team'}
            }
        }
    }
    & & \\ \\

    \theutterance \stepcounter{utterance}  
    & & & \multicolumn{2}{p{0.3\linewidth}}{
        \cellcolor[rgb]{0.9,0.9,0.9}{
            \makecell[{{p{\linewidth}}}]{
                \texttt{\tiny{[GM$|$GM]}}
                {'word': 'pivo', 'assignment': 'team'}
            }
        }
    }
    & & \\ \\

    \theutterance \stepcounter{utterance}  
    & & \multicolumn{4}{p{0.6\linewidth}}{
        \cellcolor[rgb]{0.9,0.9,0.9}{
            \makecell[{{p{\linewidth}}}]{
                \texttt{\tiny{[GM$\rangle$P2]}}
                \texttt{Beseda črv je bila uganjena, vendar je nasprotnik beseda. Vaša poteza se je tam končala.} \\
\texttt{Zdaj navedite še en seznam z vejico ločenih besed, najmanj 1 in največ 2 besed z naslednjega seznama besed, ki se najbolj nanašajo ali so najbolj povezane z besedo 'glasba'. Ne pozabite začeti odgovora z UGIBANJE:. Opozorilo: nekatere besede so bile v primerjavi s prejšnjimi zahtevami odstranjene s seznama.} \\
\\ 
\texttt{malarija, norveška, robin, kardinal, viski, sladkorna bolezen, metlica, bež, poljska, rak, oriole, mehika, stonoga, zajemalka, palec, rum, pivo, centimeter, teal, molj, hrošč, cimbala, rog} \\
            }
        }
    }
    & \\ \\

    \theutterance \stepcounter{utterance}  
    & & & \multicolumn{4}{p{0.6\linewidth}}{
        \cellcolor[rgb]{0.9,0.9,0.9}{
            \makecell[{{p{\linewidth}}}]{
                \texttt{\tiny{[GM$\langle$P2]}}
                \texttt{UGIBANJE: cimbala, rog} \\
            }
        }
    }
     \\ \\

    \theutterance \stepcounter{utterance}  
    & & & \multicolumn{2}{p{0.3\linewidth}}{
        \cellcolor[rgb]{0.9,0.9,0.9}{
            \makecell[{{p{\linewidth}}}]{
                \texttt{\tiny{[GM$|$GM]}}
                ['cimbala', 'rog']
            }
        }
    }
    & & \\ \\

    \theutterance \stepcounter{utterance}  
    & & & \multicolumn{2}{p{0.3\linewidth}}{
        \cellcolor[rgb]{0.9,0.9,0.9}{
            \makecell[{{p{\linewidth}}}]{
                \texttt{\tiny{[GM$|$GM]}}
                {'word': 'cimbala', 'assignment': 'team'}
            }
        }
    }
    & & \\ \\

    \theutterance \stepcounter{utterance}  
    & & & \multicolumn{2}{p{0.3\linewidth}}{
        \cellcolor[rgb]{0.9,0.9,0.9}{
            \makecell[{{p{\linewidth}}}]{
                \texttt{\tiny{[GM$|$GM]}}
                {'word': 'cimbala', 'assignment': 'team'}
            }
        }
    }
    & & \\ \\

    \theutterance \stepcounter{utterance}  
    & & & \multicolumn{2}{p{0.3\linewidth}}{
        \cellcolor[rgb]{0.9,0.9,0.9}{
            \makecell[{{p{\linewidth}}}]{
                \texttt{\tiny{[GM$|$GM]}}
                {'word': 'rog', 'assignment': 'opponent'}
            }
        }
    }
    & & \\ \\

    \theutterance \stepcounter{utterance}  
    & & & \multicolumn{2}{p{0.3\linewidth}}{
        \cellcolor[rgb]{0.9,0.9,0.9}{
            \makecell[{{p{\linewidth}}}]{
                \texttt{\tiny{[GM$|$GM]}}
                \texttt{rog} \\
            }
        }
    }
    & & \\ \\

    \theutterance \stepcounter{utterance}  
    & & & \multicolumn{2}{p{0.3\linewidth}}{
        \cellcolor[rgb]{0.9,0.9,0.9}{
            \makecell[{{p{\linewidth}}}]{
                \texttt{\tiny{[GM$|$GM]}}
                {'word': 'rak', 'assignment': 'opponent'}
            }
        }
    }
    & & \\ \\

    \theutterance \stepcounter{utterance}  
    & & & \multicolumn{2}{p{0.3\linewidth}}{
        \cellcolor[rgb]{0.9,0.9,0.9}{
            \makecell[{{p{\linewidth}}}]{
                \texttt{\tiny{[GM$|$GM]}}
                {'hidden': {'team': ['robin', 'palec', 'stonoga', 'malarija', 'zajemalka', 'pivo', 'bež', 'poljska'], 'innocent': ['molj', 'norveška', 'centimeter', 'metlica', 'hrošč', 'oriole', 'viski'], 'opponent': ['mehika', 'kardinal', 'sladkorna bolezen', 'teal'], 'assassin': ['rum']}, 'revealed': {'team': {'team': ['cimbala'], 'innocent': [], 'opponent': ['črv', 'rog'], 'assassin': []}, 'opponent': {'team': [], 'innocent': [], 'opponent': ['milja', 'rak'], 'assassin': []}}}
            }
        }
    }
    & & \\ \\

    \theutterance \stepcounter{utterance}  
    & & \multicolumn{4}{p{0.6\linewidth}}{
        \cellcolor[rgb]{0.9,0.9,0.9}{
            \makecell[{{p{\linewidth}}}]{
                \texttt{\tiny{[P1$\langle$GM]}}
                \texttt{Beseda cimbala je bila pravilno uganjena. Beseda rog je bila uganjena, vendar je nasprotnik beseda. Potezo vašega soigralca se je tam končala.} \\
\texttt{Sedaj podajte še en namig, povezan z nekaterimi vašimi preostalimi besedami ekipe, in seznam povezanih ciljnih besed. Ne pozabite začeti namiga z 'NAMIG: ', nato dodati novo vrstico in začeti seznam ciljnih besed, ločenih z vejico, z 'CILJI: '. Opomba: nekatere besede so bile s seznamov odstranjene v primerjavi s prejšnjimi zahtevami.} \\
\\ 
\texttt{Vaše preostale besede ekipe so: robin, palec, stonoga, malarija, zajemalka, pivo, bež, poljska.} \\
\texttt{Preostale besede nasprotnika so: rak, mehika, kardinal, sladkorna bolezen, teal.} \\
\texttt{Preostale besede, ki motijo, so: molj, norveška, centimeter, metlica, hrošč, oriole, viski.} \\
\texttt{Preostale smrtonosne besede so: rum.} \\
            }
        }
    }
    & \\ \\

    \theutterance \stepcounter{utterance}  
    & \multicolumn{4}{p{0.6\linewidth}}{
        \cellcolor[rgb]{0.9,0.9,0.9}{
            \makecell[{{p{\linewidth}}}]{
                \texttt{\tiny{[P1$\rangle$GM]}}
                \texttt{NAMIG: pijača} \\
\texttt{CILJI: pivo, zajemalka} \\
            }
        }
    }
    & & \\ \\

    \theutterance \stepcounter{utterance}  
    & & & \multicolumn{2}{p{0.3\linewidth}}{
        \cellcolor[rgb]{0.9,0.9,0.9}{
            \makecell[{{p{\linewidth}}}]{
                \texttt{\tiny{[GM$|$GM]}}
                \texttt{pijača} \\
            }
        }
    }
    & & \\ \\

    \theutterance \stepcounter{utterance}  
    & & & \multicolumn{2}{p{0.3\linewidth}}{
        \cellcolor[rgb]{0.9,0.9,0.9}{
            \makecell[{{p{\linewidth}}}]{
                \texttt{\tiny{[GM$|$GM]}}
                ['pivo', 'zajemalka']
            }
        }
    }
    & & \\ \\

    \theutterance \stepcounter{utterance}  
    & & & \multicolumn{2}{p{0.3\linewidth}}{
        \cellcolor[rgb]{0.9,0.9,0.9}{
            \makecell[{{p{\linewidth}}}]{
                \texttt{\tiny{[GM$|$GM]}}
                {'word': 'pivo', 'assignment': 'team'}
            }
        }
    }
    & & \\ \\

    \theutterance \stepcounter{utterance}  
    & & & \multicolumn{2}{p{0.3\linewidth}}{
        \cellcolor[rgb]{0.9,0.9,0.9}{
            \makecell[{{p{\linewidth}}}]{
                \texttt{\tiny{[GM$|$GM]}}
                {'word': 'zajemalka', 'assignment': 'team'}
            }
        }
    }
    & & \\ \\

    \theutterance \stepcounter{utterance}  
    & & \multicolumn{4}{p{0.6\linewidth}}{
        \cellcolor[rgb]{0.9,0.9,0.9}{
            \makecell[{{p{\linewidth}}}]{
                \texttt{\tiny{[GM$\rangle$P2]}}
                \texttt{Beseda cimbala je bila pravilno uganjena. Beseda rog je bila uganjena, vendar je nasprotnik beseda. Vaša poteza se je tam končala.} \\
\texttt{Zdaj navedite še en seznam z vejico ločenih besed, najmanj 1 in največ 2 besed z naslednjega seznama besed, ki se najbolj nanašajo ali so najbolj povezane z besedo 'pijača'. Ne pozabite začeti odgovora z UGIBANJE:. Opozorilo: nekatere besede so bile v primerjavi s prejšnjimi zahtevami odstranjene s seznama.} \\
\\ 
\texttt{malarija, norveška, robin, kardinal, viski, sladkorna bolezen, metlica, bež, poljska, oriole, mehika, stonoga, zajemalka, palec, rum, pivo, centimeter, teal, molj, hrošč} \\
            }
        }
    }
    & \\ \\

    \theutterance \stepcounter{utterance}  
    & & & \multicolumn{4}{p{0.6\linewidth}}{
        \cellcolor[rgb]{0.9,0.9,0.9}{
            \makecell[{{p{\linewidth}}}]{
                \texttt{\tiny{[GM$\langle$P2]}}
                \texttt{UGIBANJE: viski, rum} \\
            }
        }
    }
     \\ \\

    \theutterance \stepcounter{utterance}  
    & & & \multicolumn{2}{p{0.3\linewidth}}{
        \cellcolor[rgb]{0.9,0.9,0.9}{
            \makecell[{{p{\linewidth}}}]{
                \texttt{\tiny{[GM$|$GM]}}
                ['viski', 'rum']
            }
        }
    }
    & & \\ \\

    \theutterance \stepcounter{utterance}  
    & & & \multicolumn{2}{p{0.3\linewidth}}{
        \cellcolor[rgb]{0.9,0.9,0.9}{
            \makecell[{{p{\linewidth}}}]{
                \texttt{\tiny{[GM$|$GM]}}
                {'word': 'viski', 'assignment': 'innocent'}
            }
        }
    }
    & & \\ \\

    \theutterance \stepcounter{utterance}  
    & & & \multicolumn{2}{p{0.3\linewidth}}{
        \cellcolor[rgb]{0.9,0.9,0.9}{
            \makecell[{{p{\linewidth}}}]{
                \texttt{\tiny{[GM$|$GM]}}
                \texttt{viski} \\
            }
        }
    }
    & & \\ \\

    \theutterance \stepcounter{utterance}  
    & & & \multicolumn{2}{p{0.3\linewidth}}{
        \cellcolor[rgb]{0.9,0.9,0.9}{
            \makecell[{{p{\linewidth}}}]{
                \texttt{\tiny{[GM$|$GM]}}
                {'word': 'sladkorna bolezen', 'assignment': 'opponent'}
            }
        }
    }
    & & \\ \\

    \theutterance \stepcounter{utterance}  
    & & & \multicolumn{2}{p{0.3\linewidth}}{
        \cellcolor[rgb]{0.9,0.9,0.9}{
            \makecell[{{p{\linewidth}}}]{
                \texttt{\tiny{[GM$|$GM]}}
                {'hidden': {'team': ['robin', 'palec', 'stonoga', 'malarija', 'zajemalka', 'pivo', 'bež', 'poljska'], 'innocent': ['molj', 'norveška', 'centimeter', 'metlica', 'hrošč', 'oriole'], 'opponent': ['mehika', 'kardinal', 'teal'], 'assassin': ['rum']}, 'revealed': {'team': {'team': ['cimbala'], 'innocent': ['viski'], 'opponent': ['črv', 'rog'], 'assassin': []}, 'opponent': {'team': [], 'innocent': [], 'opponent': ['milja', 'rak', 'sladkorna bolezen'], 'assassin': []}}}
            }
        }
    }
    & & \\ \\

    \theutterance \stepcounter{utterance}  
    & & \multicolumn{4}{p{0.6\linewidth}}{
        \cellcolor[rgb]{0.9,0.9,0.9}{
            \makecell[{{p{\linewidth}}}]{
                \texttt{\tiny{[P1$\langle$GM]}}
                \texttt{Beseda viski je bila uganjena, vendar je nevinen beseda. Potezo vašega soigralca se je tam končala.} \\
\texttt{Sedaj podajte še en namig, povezan z nekaterimi vašimi preostalimi besedami ekipe, in seznam povezanih ciljnih besed. Ne pozabite začeti namiga z 'NAMIG: ', nato dodati novo vrstico in začeti seznam ciljnih besed, ločenih z vejico, z 'CILJI: '. Opomba: nekatere besede so bile s seznamov odstranjene v primerjavi s prejšnjimi zahtevami.} \\
\\ 
\texttt{Vaše preostale besede ekipe so: robin, palec, stonoga, malarija, zajemalka, pivo, bež, poljska.} \\
\texttt{Preostale besede nasprotnika so: mehika, kardinal, sladkorna bolezen, teal.} \\
\texttt{Preostale besede, ki motijo, so: molj, norveška, centimeter, metlica, hrošč, oriole.} \\
\texttt{Preostale smrtonosne besede so: rum.} \\
            }
        }
    }
    & \\ \\

    \theutterance \stepcounter{utterance}  
    & \multicolumn{4}{p{0.6\linewidth}}{
        \cellcolor[rgb]{0.9,0.9,0.9}{
            \makecell[{{p{\linewidth}}}]{
                \texttt{\tiny{[P1$\rangle$GM]}}
                \texttt{NAMIG: žival} \\
\texttt{CILJI: robin, stonoga} \\
            }
        }
    }
    & & \\ \\

    \theutterance \stepcounter{utterance}  
    & & & \multicolumn{2}{p{0.3\linewidth}}{
        \cellcolor[rgb]{0.9,0.9,0.9}{
            \makecell[{{p{\linewidth}}}]{
                \texttt{\tiny{[GM$|$GM]}}
                \texttt{žival} \\
            }
        }
    }
    & & \\ \\

    \theutterance \stepcounter{utterance}  
    & & & \multicolumn{2}{p{0.3\linewidth}}{
        \cellcolor[rgb]{0.9,0.9,0.9}{
            \makecell[{{p{\linewidth}}}]{
                \texttt{\tiny{[GM$|$GM]}}
                ['robin', 'stonoga']
            }
        }
    }
    & & \\ \\

    \theutterance \stepcounter{utterance}  
    & & & \multicolumn{2}{p{0.3\linewidth}}{
        \cellcolor[rgb]{0.9,0.9,0.9}{
            \makecell[{{p{\linewidth}}}]{
                \texttt{\tiny{[GM$|$GM]}}
                {'word': 'robin', 'assignment': 'team'}
            }
        }
    }
    & & \\ \\

    \theutterance \stepcounter{utterance}  
    & & & \multicolumn{2}{p{0.3\linewidth}}{
        \cellcolor[rgb]{0.9,0.9,0.9}{
            \makecell[{{p{\linewidth}}}]{
                \texttt{\tiny{[GM$|$GM]}}
                {'word': 'stonoga', 'assignment': 'team'}
            }
        }
    }
    & & \\ \\

    \theutterance \stepcounter{utterance}  
    & & \multicolumn{4}{p{0.6\linewidth}}{
        \cellcolor[rgb]{0.9,0.9,0.9}{
            \makecell[{{p{\linewidth}}}]{
                \texttt{\tiny{[GM$\rangle$P2]}}
                \texttt{Beseda viski je bila uganjena, vendar je nevinen beseda. Vaša poteza se je tam končala.} \\
\texttt{Zdaj navedite še en seznam z vejico ločenih besed, najmanj 1 in največ 2 besed z naslednjega seznama besed, ki se najbolj nanašajo ali so najbolj povezane z besedo 'žival'. Ne pozabite začeti odgovora z UGIBANJE:. Opozorilo: nekatere besede so bile v primerjavi s prejšnjimi zahtevami odstranjene s seznama.} \\
\\ 
\texttt{malarija, norveška, robin, kardinal, metlica, bež, poljska, oriole, mehika, stonoga, zajemalka, palec, rum, pivo, centimeter, teal, molj, hrošč} \\
            }
        }
    }
    & \\ \\

    \theutterance \stepcounter{utterance}  
    & & & \multicolumn{4}{p{0.6\linewidth}}{
        \cellcolor[rgb]{0.9,0.9,0.9}{
            \makecell[{{p{\linewidth}}}]{
                \texttt{\tiny{[GM$\langle$P2]}}
                \texttt{UGIBANJE: stonoga, hrošč} \\
            }
        }
    }
     \\ \\

    \theutterance \stepcounter{utterance}  
    & & & \multicolumn{2}{p{0.3\linewidth}}{
        \cellcolor[rgb]{0.9,0.9,0.9}{
            \makecell[{{p{\linewidth}}}]{
                \texttt{\tiny{[GM$|$GM]}}
                ['stonoga', 'hrošč']
            }
        }
    }
    & & \\ \\

    \theutterance \stepcounter{utterance}  
    & & & \multicolumn{2}{p{0.3\linewidth}}{
        \cellcolor[rgb]{0.9,0.9,0.9}{
            \makecell[{{p{\linewidth}}}]{
                \texttt{\tiny{[GM$|$GM]}}
                {'word': 'stonoga', 'assignment': 'team'}
            }
        }
    }
    & & \\ \\

    \theutterance \stepcounter{utterance}  
    & & & \multicolumn{2}{p{0.3\linewidth}}{
        \cellcolor[rgb]{0.9,0.9,0.9}{
            \makecell[{{p{\linewidth}}}]{
                \texttt{\tiny{[GM$|$GM]}}
                {'word': 'stonoga', 'assignment': 'team'}
            }
        }
    }
    & & \\ \\

    \theutterance \stepcounter{utterance}  
    & & & \multicolumn{2}{p{0.3\linewidth}}{
        \cellcolor[rgb]{0.9,0.9,0.9}{
            \makecell[{{p{\linewidth}}}]{
                \texttt{\tiny{[GM$|$GM]}}
                {'word': 'hrošč', 'assignment': 'innocent'}
            }
        }
    }
    & & \\ \\

    \theutterance \stepcounter{utterance}  
    & & & \multicolumn{2}{p{0.3\linewidth}}{
        \cellcolor[rgb]{0.9,0.9,0.9}{
            \makecell[{{p{\linewidth}}}]{
                \texttt{\tiny{[GM$|$GM]}}
                \texttt{hrošč} \\
            }
        }
    }
    & & \\ \\

    \theutterance \stepcounter{utterance}  
    & & & \multicolumn{2}{p{0.3\linewidth}}{
        \cellcolor[rgb]{0.9,0.9,0.9}{
            \makecell[{{p{\linewidth}}}]{
                \texttt{\tiny{[GM$|$GM]}}
                {'word': 'teal', 'assignment': 'opponent'}
            }
        }
    }
    & & \\ \\

    \theutterance \stepcounter{utterance}  
    & & & \multicolumn{2}{p{0.3\linewidth}}{
        \cellcolor[rgb]{0.9,0.9,0.9}{
            \makecell[{{p{\linewidth}}}]{
                \texttt{\tiny{[GM$|$GM]}}
                {'hidden': {'team': ['robin', 'palec', 'malarija', 'zajemalka', 'pivo', 'bež', 'poljska'], 'innocent': ['molj', 'norveška', 'centimeter', 'metlica', 'oriole'], 'opponent': ['mehika', 'kardinal'], 'assassin': ['rum']}, 'revealed': {'team': {'team': ['cimbala', 'stonoga'], 'innocent': ['viski', 'hrošč'], 'opponent': ['črv', 'rog'], 'assassin': []}, 'opponent': {'team': [], 'innocent': [], 'opponent': ['milja', 'rak', 'sladkorna bolezen', 'teal'], 'assassin': []}}}
            }
        }
    }
    & & \\ \\

    \theutterance \stepcounter{utterance}  
    & & \multicolumn{4}{p{0.6\linewidth}}{
        \cellcolor[rgb]{0.9,0.9,0.9}{
            \makecell[{{p{\linewidth}}}]{
                \texttt{\tiny{[P1$\langle$GM]}}
                \texttt{Beseda stonoga je bila pravilno uganjena. Beseda hrošč je bila uganjena, vendar je nevinen beseda. Potezo vašega soigralca se je tam končala.} \\
\texttt{Sedaj podajte še en namig, povezan z nekaterimi vašimi preostalimi besedami ekipe, in seznam povezanih ciljnih besed. Ne pozabite začeti namiga z 'NAMIG: ', nato dodati novo vrstico in začeti seznam ciljnih besed, ločenih z vejico, z 'CILJI: '. Opomba: nekatere besede so bile s seznamov odstranjene v primerjavi s prejšnjimi zahtevami.} \\
\\ 
\texttt{Vaše preostale besede ekipe so: robin, palec, malarija, zajemalka, pivo, bež, poljska.} \\
\texttt{Preostale besede nasprotnika so: mehika, kardinal, teal.} \\
\texttt{Preostale besede, ki motijo, so: molj, norveška, centimeter, metlica, oriole.} \\
\texttt{Preostale smrtonosne besede so: rum.} \\
            }
        }
    }
    & \\ \\

    \theutterance \stepcounter{utterance}  
    & \multicolumn{4}{p{0.6\linewidth}}{
        \cellcolor[rgb]{0.9,0.9,0.9}{
            \makecell[{{p{\linewidth}}}]{
                \texttt{\tiny{[P1$\rangle$GM]}}
                \texttt{NAMIG: ptica} \\
\texttt{CILJI: robin} \\
            }
        }
    }
    & & \\ \\

    \theutterance \stepcounter{utterance}  
    & & & \multicolumn{2}{p{0.3\linewidth}}{
        \cellcolor[rgb]{0.9,0.9,0.9}{
            \makecell[{{p{\linewidth}}}]{
                \texttt{\tiny{[GM$|$GM]}}
                \texttt{ptica} \\
            }
        }
    }
    & & \\ \\

    \theutterance \stepcounter{utterance}  
    & & & \multicolumn{2}{p{0.3\linewidth}}{
        \cellcolor[rgb]{0.9,0.9,0.9}{
            \makecell[{{p{\linewidth}}}]{
                \texttt{\tiny{[GM$|$GM]}}
                ['robin']
            }
        }
    }
    & & \\ \\

    \theutterance \stepcounter{utterance}  
    & & & \multicolumn{2}{p{0.3\linewidth}}{
        \cellcolor[rgb]{0.9,0.9,0.9}{
            \makecell[{{p{\linewidth}}}]{
                \texttt{\tiny{[GM$|$GM]}}
                {'word': 'robin', 'assignment': 'team'}
            }
        }
    }
    & & \\ \\

    \theutterance \stepcounter{utterance}  
    & & \multicolumn{4}{p{0.6\linewidth}}{
        \cellcolor[rgb]{0.9,0.9,0.9}{
            \makecell[{{p{\linewidth}}}]{
                \texttt{\tiny{[GM$\rangle$P2]}}
                \texttt{Beseda stonoga je bila pravilno uganjena. Beseda hrošč je bila uganjena, vendar je nevinen beseda. Vaša poteza se je tam končala.} \\
\texttt{Zdaj navedite še en seznam z vejico ločenih besed, najmanj 1 in največ 1 besed z naslednjega seznama besed, ki se najbolj nanašajo ali so najbolj povezane z besedo 'ptica'. Ne pozabite začeti odgovora z UGIBANJE:. Opozorilo: nekatere besede so bile v primerjavi s prejšnjimi zahtevami odstranjene s seznama.} \\
\\ 
\texttt{malarija, norveška, robin, kardinal, metlica, bež, poljska, oriole, mehika, zajemalka, palec, rum, pivo, centimeter, molj} \\
            }
        }
    }
    & \\ \\

    \theutterance \stepcounter{utterance}  
    & & & \multicolumn{4}{p{0.6\linewidth}}{
        \cellcolor[rgb]{0.9,0.9,0.9}{
            \makecell[{{p{\linewidth}}}]{
                \texttt{\tiny{[GM$\langle$P2]}}
                \texttt{UGIBANJE: robin} \\
            }
        }
    }
     \\ \\

    \theutterance \stepcounter{utterance}  
    & & & \multicolumn{2}{p{0.3\linewidth}}{
        \cellcolor[rgb]{0.9,0.9,0.9}{
            \makecell[{{p{\linewidth}}}]{
                \texttt{\tiny{[GM$|$GM]}}
                ['robin']
            }
        }
    }
    & & \\ \\

    \theutterance \stepcounter{utterance}  
    & & & \multicolumn{2}{p{0.3\linewidth}}{
        \cellcolor[rgb]{0.9,0.9,0.9}{
            \makecell[{{p{\linewidth}}}]{
                \texttt{\tiny{[GM$|$GM]}}
                {'word': 'robin', 'assignment': 'team'}
            }
        }
    }
    & & \\ \\

    \theutterance \stepcounter{utterance}  
    & & & \multicolumn{2}{p{0.3\linewidth}}{
        \cellcolor[rgb]{0.9,0.9,0.9}{
            \makecell[{{p{\linewidth}}}]{
                \texttt{\tiny{[GM$|$GM]}}
                {'word': 'robin', 'assignment': 'team'}
            }
        }
    }
    & & \\ \\

    \theutterance \stepcounter{utterance}  
    & & & \multicolumn{2}{p{0.3\linewidth}}{
        \cellcolor[rgb]{0.9,0.9,0.9}{
            \makecell[{{p{\linewidth}}}]{
                \texttt{\tiny{[GM$|$GM]}}
                {'word': 'mehika', 'assignment': 'opponent'}
            }
        }
    }
    & & \\ \\

    \theutterance \stepcounter{utterance}  
    & & & \multicolumn{2}{p{0.3\linewidth}}{
        \cellcolor[rgb]{0.9,0.9,0.9}{
            \makecell[{{p{\linewidth}}}]{
                \texttt{\tiny{[GM$|$GM]}}
                {'hidden': {'team': ['palec', 'malarija', 'zajemalka', 'pivo', 'bež', 'poljska'], 'innocent': ['molj', 'norveška', 'centimeter', 'metlica', 'oriole'], 'opponent': ['kardinal'], 'assassin': ['rum']}, 'revealed': {'team': {'team': ['cimbala', 'stonoga', 'robin'], 'innocent': ['viski', 'hrošč'], 'opponent': ['črv', 'rog'], 'assassin': []}, 'opponent': {'team': [], 'innocent': [], 'opponent': ['milja', 'rak', 'sladkorna bolezen', 'teal', 'mehika'], 'assassin': []}}}
            }
        }
    }
    & & \\ \\

    \theutterance \stepcounter{utterance}  
    & & \multicolumn{4}{p{0.6\linewidth}}{
        \cellcolor[rgb]{0.9,0.9,0.9}{
            \makecell[{{p{\linewidth}}}]{
                \texttt{\tiny{[P1$\langle$GM]}}
                \texttt{Beseda robin je bila pravilno uganjena. Potezo vašega soigralca se je tam končala.} \\
\texttt{Sedaj podajte še en namig, povezan z nekaterimi vašimi preostalimi besedami ekipe, in seznam povezanih ciljnih besed. Ne pozabite začeti namiga z 'NAMIG: ', nato dodati novo vrstico in začeti seznam ciljnih besed, ločenih z vejico, z 'CILJI: '. Opomba: nekatere besede so bile s seznamov odstranjene v primerjavi s prejšnjimi zahtevami.} \\
\\ 
\texttt{Vaše preostale besede ekipe so: palec, malarija, zajemalka, pivo, bež, poljska.} \\
\texttt{Preostale besede nasprotnika so: mehika, kardinal.} \\
\texttt{Preostale besede, ki motijo, so: molj, norveška, centimeter, metlica, oriole.} \\
\texttt{Preostale smrtonosne besede so: rum.} \\
            }
        }
    }
    & \\ \\

    \theutterance \stepcounter{utterance}  
    & \multicolumn{4}{p{0.6\linewidth}}{
        \cellcolor[rgb]{0.9,0.9,0.9}{
            \makecell[{{p{\linewidth}}}]{
                \texttt{\tiny{[P1$\rangle$GM]}}
                \texttt{NAMIG: dež} \\
\texttt{CILJI: malarija, pivo} \\
            }
        }
    }
    & & \\ \\

    \theutterance \stepcounter{utterance}  
    & & & \multicolumn{2}{p{0.3\linewidth}}{
        \cellcolor[rgb]{0.9,0.9,0.9}{
            \makecell[{{p{\linewidth}}}]{
                \texttt{\tiny{[GM$|$GM]}}
                \texttt{dež} \\
            }
        }
    }
    & & \\ \\

    \theutterance \stepcounter{utterance}  
    & & & \multicolumn{2}{p{0.3\linewidth}}{
        \cellcolor[rgb]{0.9,0.9,0.9}{
            \makecell[{{p{\linewidth}}}]{
                \texttt{\tiny{[GM$|$GM]}}
                ['malarija', 'pivo']
            }
        }
    }
    & & \\ \\

    \theutterance \stepcounter{utterance}  
    & & & \multicolumn{2}{p{0.3\linewidth}}{
        \cellcolor[rgb]{0.9,0.9,0.9}{
            \makecell[{{p{\linewidth}}}]{
                \texttt{\tiny{[GM$|$GM]}}
                {'word': 'malarija', 'assignment': 'team'}
            }
        }
    }
    & & \\ \\

    \theutterance \stepcounter{utterance}  
    & & & \multicolumn{2}{p{0.3\linewidth}}{
        \cellcolor[rgb]{0.9,0.9,0.9}{
            \makecell[{{p{\linewidth}}}]{
                \texttt{\tiny{[GM$|$GM]}}
                {'word': 'pivo', 'assignment': 'team'}
            }
        }
    }
    & & \\ \\

    \theutterance \stepcounter{utterance}  
    & & \multicolumn{4}{p{0.6\linewidth}}{
        \cellcolor[rgb]{0.9,0.9,0.9}{
            \makecell[{{p{\linewidth}}}]{
                \texttt{\tiny{[GM$\rangle$P2]}}
                \texttt{Beseda robin je bila pravilno uganjena. Vaša poteza se je tam končala.} \\
\texttt{Zdaj navedite še en seznam z vejico ločenih besed, najmanj 1 in največ 2 besed z naslednjega seznama besed, ki se najbolj nanašajo ali so najbolj povezane z besedo 'dež'. Ne pozabite začeti odgovora z UGIBANJE:. Opozorilo: nekatere besede so bile v primerjavi s prejšnjimi zahtevami odstranjene s seznama.} \\
\\ 
\texttt{malarija, norveška, kardinal, metlica, bež, poljska, oriole, zajemalka, palec, rum, pivo, centimeter, molj} \\
            }
        }
    }
    & \\ \\

    \theutterance \stepcounter{utterance}  
    & & & \multicolumn{4}{p{0.6\linewidth}}{
        \cellcolor[rgb]{0.9,0.9,0.9}{
            \makecell[{{p{\linewidth}}}]{
                \texttt{\tiny{[GM$\langle$P2]}}
                \texttt{UGIBANJE: poljska} \\
            }
        }
    }
     \\ \\

    \theutterance \stepcounter{utterance}  
    & & & \multicolumn{2}{p{0.3\linewidth}}{
        \cellcolor[rgb]{0.9,0.9,0.9}{
            \makecell[{{p{\linewidth}}}]{
                \texttt{\tiny{[GM$|$GM]}}
                ['poljska']
            }
        }
    }
    & & \\ \\

    \theutterance \stepcounter{utterance}  
    & & & \multicolumn{2}{p{0.3\linewidth}}{
        \cellcolor[rgb]{0.9,0.9,0.9}{
            \makecell[{{p{\linewidth}}}]{
                \texttt{\tiny{[GM$|$GM]}}
                {'word': 'poljska', 'assignment': 'team'}
            }
        }
    }
    & & \\ \\

    \theutterance \stepcounter{utterance}  
    & & & \multicolumn{2}{p{0.3\linewidth}}{
        \cellcolor[rgb]{0.9,0.9,0.9}{
            \makecell[{{p{\linewidth}}}]{
                \texttt{\tiny{[GM$|$GM]}}
                {'word': 'kardinal', 'assignment': 'opponent'}
            }
        }
    }
    & & \\ \\

    \theutterance \stepcounter{utterance}  
    & & & \multicolumn{2}{p{0.3\linewidth}}{
        \cellcolor[rgb]{0.9,0.9,0.9}{
            \makecell[{{p{\linewidth}}}]{
                \texttt{\tiny{[GM$|$GM]}}
                {'hidden': {'team': ['palec', 'malarija', 'zajemalka', 'pivo', 'bež'], 'innocent': ['molj', 'norveška', 'centimeter', 'metlica', 'oriole'], 'opponent': [], 'assassin': ['rum']}, 'revealed': {'team': {'team': ['cimbala', 'stonoga', 'robin', 'poljska'], 'innocent': ['viski', 'hrošč'], 'opponent': ['črv', 'rog'], 'assassin': []}, 'opponent': {'team': [], 'innocent': [], 'opponent': ['milja', 'rak', 'sladkorna bolezen', 'teal', 'mehika', 'kardinal'], 'assassin': []}}}
            }
        }
    }
    & & \\ \\

    \theutterance \stepcounter{utterance}  
    & & \multicolumn{4}{p{0.6\linewidth}}{
        \cellcolor[rgb]{0.9,0.9,0.9}{
            \makecell[{{p{\linewidth}}}]{
                \texttt{\tiny{[P1$\langle$GM]}}
                \texttt{Beseda poljska je bila pravilno uganjena. Potezo vašega soigralca se je tam končala.} \\
\texttt{Sedaj podajte še en namig, povezan z nekaterimi vašimi preostalimi besedami ekipe, in seznam povezanih ciljnih besed. Ne pozabite začeti namiga z 'NAMIG: ', nato dodati novo vrstico in začeti seznam ciljnih besed, ločenih z vejico, z 'CILJI: '. Opomba: nekatere besede so bile s seznamov odstranjene v primerjavi s prejšnjimi zahtevami.} \\
\\ 
\texttt{Vaše preostale besede ekipe so: palec, malarija, zajemalka, pivo, bež.} \\
\texttt{Preostale besede nasprotnika so: kardinal.} \\
\texttt{Preostale besede, ki motijo, so: molj, norveška, centimeter, metlica, oriole.} \\
\texttt{Preostale smrtonosne besede so: rum.} \\
            }
        }
    }
    & \\ \\

    \theutterance \stepcounter{utterance}  
    & \multicolumn{4}{p{0.6\linewidth}}{
        \cellcolor[rgb]{0.9,0.9,0.9}{
            \makecell[{{p{\linewidth}}}]{
                \texttt{\tiny{[P1$\rangle$GM]}}
                \texttt{NAMIG: prst} \\
\texttt{CILJI: palec, zajemalka} \\
            }
        }
    }
    & & \\ \\

    \theutterance \stepcounter{utterance}  
    & & & \multicolumn{2}{p{0.3\linewidth}}{
        \cellcolor[rgb]{0.9,0.9,0.9}{
            \makecell[{{p{\linewidth}}}]{
                \texttt{\tiny{[GM$|$GM]}}
                \texttt{prst} \\
            }
        }
    }
    & & \\ \\

    \theutterance \stepcounter{utterance}  
    & & & \multicolumn{2}{p{0.3\linewidth}}{
        \cellcolor[rgb]{0.9,0.9,0.9}{
            \makecell[{{p{\linewidth}}}]{
                \texttt{\tiny{[GM$|$GM]}}
                ['palec', 'zajemalka']
            }
        }
    }
    & & \\ \\

    \theutterance \stepcounter{utterance}  
    & & & \multicolumn{2}{p{0.3\linewidth}}{
        \cellcolor[rgb]{0.9,0.9,0.9}{
            \makecell[{{p{\linewidth}}}]{
                \texttt{\tiny{[GM$|$GM]}}
                {'word': 'palec', 'assignment': 'team'}
            }
        }
    }
    & & \\ \\

    \theutterance \stepcounter{utterance}  
    & & & \multicolumn{2}{p{0.3\linewidth}}{
        \cellcolor[rgb]{0.9,0.9,0.9}{
            \makecell[{{p{\linewidth}}}]{
                \texttt{\tiny{[GM$|$GM]}}
                {'word': 'zajemalka', 'assignment': 'team'}
            }
        }
    }
    & & \\ \\

    \theutterance \stepcounter{utterance}  
    & & & \multicolumn{2}{p{0.3\linewidth}}{
        \cellcolor[rgb]{0.9,0.9,0.9}{
            \makecell[{{p{\linewidth}}}]{
                \texttt{\tiny{[GM$|$GM]}}
                \texttt{opponent has won} \\
            }
        }
    }
    & & \\ \\

\end{supertabular}
}

\end{document}
